% ═══════════════════════════════════════════════════════════════════════════════
%  INTRODUCTION GÉNÉRALE
% ═══════════════════════════════════════════════════════════════════════════════
\newpage
\section*{Introduction Générale}
\addcontentsline{toc}{section}{Introduction Générale}
\label{sec:introduction_generale}

\subsection*{Contexte général}
\label{intro:contexte}
\addcontentsline{toc}{subsection}{Contexte général}

L'organisation des soutenances en milieu universitaire est une activité critique qui se heurte aujourd'hui à plusieurs défis. La multiplicité des acteurs -- étudiants, encadrants, membres de jury et services administratifs -- nécessite une coordination fine et rigoureuse. Pourtant, les processus traditionnels reposent encore largement sur des tableurs Excel, des échanges d'e-mails massifs et des documents papier, engendrant une perte d'efficacité significative et des erreurs de planification fréquentes.

Dans ce contexte, la faculté des Sciences Ben M'Sick (FSBM) de l'Université Hassan II de Casablanca, à travers son département de Mathématiques, Informatique et Physique (MIP), cherche à moderniser ses outils de gestion académique. L'objectif est de centraliser les données et d'automatiser les flux de travail via une plateforme web dédiée, capable de répondre aux exigences croissantes de qualité et de réactivité.

\subsection*{Problématique}
\label{intro:problematique}
\addcontentsline{toc}{subsection}{Problématique}

La coordination des soutenances dans un département universitaire fait face à une complexité croissante. Les contraintes à respecter sont multiples et souvent interdépendantes. La disponibilité des enseignants, qui interviennent à la fois comme encadrants et comme membres de jury, doit être prise en compte pour éviter les doubles affectations horaires. La réservation des salles, dont les capacités varient, nécessite une allocation cohérente avec le nombre de participants attendus.

Par ailleurs, la production des documents administratifs accompagnant chaque soutenance -- convocations, fiches d'évaluation -- représente une charge de travail récurrente et chronophage. La saisie et le suivi des évaluations, la transmission des notes aux services compétents et la génération des rapports de synthèse complètent ce tableau.

Le cœur de la problématique réside donc dans l'absence d'un outil unifié capable de gérer l'ensemble de ces aspects de manière intégrée et automatisée. L'organisation manuelle actuelle, basée sur des outils génériques non adaptés à cette finalité spécifique, génère des inefficacités, des erreurs et une traçabilité insuffisante.

\subsection*{Objectifs du projet}
\label{intro:objectifs}
\addcontentsline{toc}{subsection}{Objectifs du projet}

\begin{projectbox}{Objectif principal}
    Développer une application web complète et fonctionnelle permettant la gestion centralisée, automatisée et fiable de l'organisation des soutenances et des jurys au sein d'un établissement universitaire.
\end{projectbox}

Cette vision générale se décline en plusieurs objectifs spécifiques. Le système doit assurer une gestion rigoureuse des acteurs académiques, incluant les étudiants, les encadrants, les membres de jury et les responsables pédagogiques. La planification des soutenances constitue un axe majeur, intégrant la gestion des calendriers, des salles, ainsi que l'affectation automatique des jurys tout en détectant les conflits de manière proactive. Par ailleurs, la production documentaire est automatisée pour la génération des convocations, fiches d'évaluation et procès-verbaux, tandis que le module d'évaluation permet une saisie fiable des notes, des observations et l'enregistrement formel des délibérations.

\subsection*{Méthodologie suivie}
\label{intro:methodologie}
\addcontentsline{toc}{subsection}{Méthodologie suivie}

La réalisation de ce projet a été menée selon une démarche structurée en quatre phases successives. Le processus a débuté par une phase d'analyse de l'existant, visant à étudier les mécanismes actuels d'organisation des soutenances, à en identifier les limites opérationnelles et à formaliser les besoins réels. Cette étape a posé les bases pour la phase de conception, consacrée à la modélisation UML—incluant les cas d'utilisation, les diagrammes de classes, de séquence et d'activité—ainsi qu'à la définition de l'architecture technique et de la modélisation de la base de données. S'en est suivie la phase de réalisation, focalisée sur l'implémentation effective du backend Spring Boot et du frontend React, pour aboutir finalement à la phase de tests et de validation, cruciale pour garantir la conformité fonctionnelle et la robustesse du système avant son déploiement.

\subsection*{Structure du rapport}
\label{intro:structure}
\addcontentsline{toc}{subsection}{Structure du rapport}

Le présent rapport est organisé en cinq chapitres distincts. Le premier chapitre présente le contexte du stage au sein de la Faculté des Sciences Ben M'Sick. Le deuxième chapitre dresse un état des lieux du système actuel, examine les solutions comparables et justifie les choix technologiques effectués. Le troisième chapitre détaille l'analyse des besoins et la conception complète du système. Le quatrième chapitre décrit la réalisation technique et les solutions apportées aux défis rencontrés. Enfin, le cinquième chapitre expose la stratégie de test et les modalités de validation, avant qu'une conclusion générale ne dresse le bilan du travail et ne trace les perspectives d'évolution.
